\documentclass[12pt]{article}
\usepackage[T1]{fontenc}
\usepackage[utf8]{inputenc}
\usepackage{lmodern}
\usepackage{textcomp}
\usepackage{enumitem}
\usepackage{geometry}
\geometry{margin=1in}
\begin{document}
\begin{center}
Answer Key - History - Undergraduate Level Assessment 13 \
\small (Answers are ordered exactly as in the question paper.)
\end{center}

\section*{Multiple Choice}
\begin{enumerate}[label=\arabic*.]
\item \textbf{Answer:} A
\\ \textit{Options:} A) bonobos; B) orangutans; C) gibbons; D) chimpanzees; E) baboons
\\ \textit{Note:} Humans share a more recent common ancestor with bonobos and chimpanzees than with any other primates, making them our closest living relatives in the evolutionary tree.
\item \textbf{Answer:} E
\\ \textit{Options:} A) Fire can be used in warfare and to clear fields.; B) Fire was used to create the first forms of currency.; C) Fire allowed early humans to fly kites at night.; D) Fire was used to create intricate cave paintings.; E) Fire makes meat more digestible and kills bacteria.
\\ \textit{Note:} The correct answer highlights that the use of fire in cooking not only enhances the digestibility of meat for early humans but also reduces the risk of foodborne illnesses by killing harmful bacteria, thereby improving nutrition and overall health, which were crucial for survival and evolutionary advancement.
\item \textbf{Answer:} C
\\ \textit{Options:} A) was a simple form of writing with one symbol representing one word; B) all of the above; C) appears to have developed suddenly; D) did not use pictographs; E) was borrowed from the Sumerians
\\ \textit{Note:} The Egyptian system of hieroglyphics is considered to have developed suddenly due to the lack of evidence for gradual evolution from simpler forms of writing, indicating a rapid establishment of a complex writing system in ancient Egypt.
\item \textbf{Answer:} A
\\ \textit{Options:} A) Europe; B) Africa; C) Central America; D) Asia; E) None of the above
\\ \textit{Note:} Europe is recognized for housing the earliest known examples of rock art, specifically etchings and engravings, such as those found in the caves of Lascaux and Chauvet, which date back to the Upper Paleolithic period.
\item \textbf{Answer:} A
\\ \textit{Options:} A) barley, lentils, and peas.; B) soybeans, emmer, and rice.; C) millet, barley, and maize.; D) wheat, rice, and peas.; E) rice, soybeans, and maize.
\\ \textit{Note:} The excavation of Yang-shao sites in China reveals evidence of early agricultural practices, specifically the domestication of barley, lentils, and peas, which were significant crops cultivated by Neolithic communities in this region.
\end{enumerate}

\section*{True / False}
\begin{enumerate}[label=\arabic*.]
\setcounter{enumi}{5}
\item \textbf{Answer:} True
\\ \textit{Options:} A) True; B) False
\\ \textit{Note:} According to the passage: Vijayanagar Empire
\item \textbf{Answer:} False
\\ \textit{Options:} A) True; B) False
\\ \textit{Note:} The correct answer is: around May 8
\item \textbf{Answer:} False
\\ \textit{Options:} A) True; B) False
\\ \textit{Note:} The correct answer is: Jon Pertwee
\item \textbf{Answer:} False
\\ \textit{Options:} A) True; B) False
\\ \textit{Note:} The correct answer is: Ángel Trías
\item \textbf{Answer:} True
\\ \textit{Options:} A) True; B) False
\\ \textit{Note:} According to the passage: Catalonia and the Basque provinces
\end{enumerate}

\section*{Long Form Response}
\begin{enumerate}[label=\arabic*.]
\setcounter{enumi}{10}
\item \textbf{Answer:} nine
\\ \textit{Note:} Expected answer: nine
\item \textbf{Answer:} 2014
\\ \textit{Note:} Expected answer: 2014
\end{enumerate}

\vfill
\noindent\textit{End of Answer Key}
\end{document}